\chapter{RESUMO}\label{chap:resumo}

Este trabalho propõe o desenvolvimento do SIRAS --- Sistema de Interpretação
e Recomendação Agronômica de Solo ---, um sistema especialista baseado em
regras para interpretação automática de análises químicas e físicas de solo,
geração de recomendações técnicas de correção e manejo, e estimativa de aptidão
agrícola para 61 culturas e grupos de culturas com relevância para o Rio Grande do Sul, abrangendo
culturas de grãos, hortaliças, tubérculos, frutíferas, erva-mate, tabaco e cana-de-açúcar.
A motivação decorre da observação, em contextos de assistência técnica rural e
crédito agrícola, de que a ausência de ferramentas computacionais acessíveis e
integradas compromete a qualidade das decisões técnicas e financeiras no setor
agropecuário. A revisão de literatura abordará os parâmetros químicos e físicos
do solo relevantes para a interpretação agronômica, os critérios de calagem e
adubação vigentes para o RS conforme o Manual de Calagem e Adubação da Sociedade
Brasileira de Ciência do Solo, os fundamentos de sistemas especialistas baseados
em regras e os sistemas computacionais correlatos existentes, como o FertFacil e
o Sistema de Avaliação de Aptidão Agrícola da Embrapa. A metodologia
caracteriza-se como pesquisa descritiva e aplicada, de abordagem quantitativa,
e prevê a estruturação da base de conhecimento a partir das normas regionais,
o projeto e a implementação dos módulos de entrada de dados, recomendação e
estimativa de aptidão, e a validação do SIRAS em três dimensões complementares:
a correção das saídas, medida pela taxa de concordância com valores de
referência calculados de forma independente, que definirá a aceitação das
hipóteses formuladas; a comparação qualitativa com ferramentas existentes; e a
avaliação de usabilidade junto ao público-alvo. Espera-se, ao final, dispor de uma ferramenta que integre recomendação
de manejo e estimativa de aptidão agrícola a partir de uma análise individual de
solo, preenchendo uma lacuna identificada nos sistemas disponíveis e contribuindo
para a tomada de decisão de técnicos, produtores e instituições financeiras que
operam com crédito rural.
